\documentclass[12pt]{article}

\usepackage[a4paper, total={6in, 8in}]{geometry}
\usepackage{textcomp}

\begin{document}
\setlength{\parindent}{0pt}

\def \t {\textrightarrow}
\def \v {\vspace{1em}}
\def \bi {\begin{itemize}}
\def \ei {\end{itemize}}
\def \s[#1] {\section*{#1}}
\def \ss[#1] {\subsection*{#1}}
\def \sss[#1] {\subsubsection*{#1}}

\s[Sartre]
Definisce come gli altri fil. l'uomo come pura possiblita \t quindi pura liberta, che genera angoscia \t uomo non ha indicatore che da direzione \\
L'angoscia è l'esperienza del nulla \t ed è l'esperienza della liberta incondizionata, liberta assoluta \\
L'altra faccia della liberta assoluta è la responsabilita \t l uomo è tutto cio che progetta di essere, e quindi mi devo assumere la responsabilita di questa scelta e delle sue conseguenze \\
Chi cerca delle scuse alla responsabilita è in malafede \t se la liberta è assoluta, anche la scelta lo è \\
Se io vengo per esempio mobilitato in guerra, se io combatto diventa la mia guerra \t oppure posso o disertare, o suicidarmi \t sono delle soluzioni estreme, ma sono delle alternative \\
La necessita esclude ogni forma di alternativa \t l'alternativa non esiste \\
Di fronte alla chiamata alle armi aleternativa esiste, posso disterarmi o uccidermi \t sono opzioni che io giudico impraticapibili, ma esistono \t non siamo nell orizzonte della necessita \\
Quando io le escludo, perche esistono, io scelgo la guerra \t e la scelgo assolutamente

\v

Liberta incondizionata = angosica \t che è aggravata dall'assoluta responsabilita \\
Liberta incondizionata dice che non posso non fare delle scelte \t che pero faccio nella totale contingenza \t alla fine non sono importanti \t che io le faccia o no, io mi sto progettando verso il nulla \t sono finito \\
Alla luce di questo la vita è un avventura assurda \t io mi continuo a progettare in una dimensione di nulla di senso \t i miei progetti non hanno significato, perche finiranno \\
L'assurdo della vita sta nel fatto che io non posso non progettarmi \t l uomo è una passione inutile \t non puo fare altro che realizzarsi in progetti in orizzonte di totale assenza di senso \\
La liberta consiste nella scelta del proprio essere, che è assurda \t perche si va verso il nulla

\v

L uomo oltre a essere coscienza è anche essere per altri pero \t e l'esistenza degli altri non deve essere dimostrata, non arriva neanche dopo la mia di esistenza \t come H. \\
Io colgo gli altri esistenti in modo immediato \t quando gli altri invadono il campo della mia soggettivita \\
Io mi rendo conto degli altri \t li colgo come altro da me, quando l'altro prova a trasformarmi da soggetto in oggetto \t quando prova a rendermi un oggetto del suo progetto \\
Il mio progetto puo coinvolgere anche gli altri uomini \t e gli altri si rendono conto di essere altro da me quando io li vedo come oggetti per il mio progetto \t cosi sto limitando o distruggendo la tua possiblita di essere quello che vuoi \\
Rapporto nasce da un conflitto \t mentre in H. prendersi cura \\
Qua invece colgo gli altri quando gli altri compiono una pressione su di me \\
Ogni rapporto umano è conflitto \t l'inferno sono gli altri, no dimensione di rapporto pacifico \\
Io mi accorgo che l'altro vuole fare di me l oggetto di un suo progetto quando l'altro mi guarda \t l'altro non è questo oggetto che io guardo, ma io mi accorgo dell altro quando l altro pone il suo sguardo su di me \\
È uno sguardo oggettivante \t mi guarda come potenziale oggetto del suo progetto \t e noi facciamo esperienza di questo con complessi di inferiorita, e quindi provo timidezza, pudore, \\
Sono messo in condizione di inferiorita \t perche mi sento guardato in modo che svilisce il mio essere soggetto \\
Quando nella mia coscienza entra un altro, la mia esperienza viene modificata \t io non sono + centro della mia esperienza, perdo progettualita perche entro come parte del progetto di un altro \t non mi sento compiuto e quindi sento inferiorita \\
Mi vergogno per esempio quando un altro mi guarda in modo giudicante \t da qui nasce il pudore \t lo sguardo che gli altri poggiano su di me scaturisce questo

\v

Quando altro entra nell orizzonte del mio progetto nasce il conflitto \t perche io mi difendo: anche io faccio dell altro oggetto del mio progetto \\
Anche il mio sguardo è oggettivante \\
Sentimenti fondamentali sono amore e odio \t i sentimenti dove questa modalita conflittuale si evidenzia di + \\
Perche nell'amore il mio progetto è quello di farmi amare da te \t fa di te un oggetto del mio progetto, vorrei anche che tu mi ami \t voglio che tu rispondi al mio progetto \\
Nell'amore faccio dell altro l'oggetto del mio progetto \t non un altro soggetto \t i progetti sono diversi, quindi comunque c'è conflitto \t pero soggettivita possono riemergere \\
Ti amo come dico io = mi devi amare nel mio progetto \t se l'altro fa lo stesso, allora scontro \t nell'amore io tendo comunque a schiacciare l'altro, voglio negare la sua liberta \\
Nell'odio invece io riconosco la liberta dell altro, e come opposta alla mia \t e quindi la voglio negare, ma non ci riesco quindi lo odio \\
Il progetto d'amore è egoista, vuole negare liberta dell altro \t quando non riesco vorrei negarla, voglio renderla oggetto \t ma non riesco quindi nasce odio \\
Quando odio mi scontro con la liberta dell'altro \\
Anche per questo uomp è pasione inutile \t creando relazioni d'amore, non uscira mai da dimensione di conflitto \t ciascuno di noi è un carnefice per gli altri \\
Sartre ha filosofia disperata

\v

Poi lui smorza questo tono cosi acceso \t nel 46 si vede tono + pacato, come in "Esistenzialismo è un umanismo" \t scritto in risposta a critiche dei marxisti riguardo alla liberta assoluta \\
Lui dice che non c'è determinismo per uomo \t è completa liberta, responsabilita è totale etc. \\
In quest opera delinea come una morale sociale \t modifica il senso del conflitto \t prima diceva che rapporti sereni con altri non si possono avere \\
Ora dice che la liberta nostra si intreccia con quelle degli altri \t noi vogliamo la nostra liberta, che pero dipende anche dal fatto che gli altri la possono esercitare \t non possiamo essere liberi solo noi \\
Liberta in definizione appartiene essenzialmente all uomo \t ma socialmente non posso difendere la mia liberta senza volere quella degli altri \t la mia liberta non puo prescindere dalla liberta altrui \\
Ciascuono deve poter esercitare liberamente la sua liberta perche lo possa fare anche io \\
Ammorbidimento si ha perche ha passata la guerra \t ha lasciato traccia \t se non ci fosse stato senso di appartenenza, di comunita di resistenza, ha influenzato il suo pensiero 

\v

La mia liberta pero non dipende solo da quella degli altri \t uomo deve fare i conti con precise situazioni, che possono genrare ostacoli a mia liberta \\
Ho bisogno di liberta degli atlri (socialmente) ma devo anche tenere conto di ostacoli \\
Affronta quindi la questione dei rapporti tra esistenzialismo e marxismo \t scrive vari saggi e scrive "Critica della ragione dialettica" (1960) \\
Qua lui esprime la sua opinione sul marxismo \t e come la sua visione marxista si possa integrare con la sua visione \\
Lui è un marxista \t lui aderisce al materialismo storico \t il modo di produzione definisce la vita dell uomo, la struttura della storia è economica \\
Lui invece rifiuta il materialismo dialettico \t perche il marxismo non puo essere mat. dialet. se con questo si intende la struttura metafisica che vuole ingabbiare la natura, idea che ci sia evoluzione metafisicad/dialettica della natura \\
Dialettica potrebbe esistere \t ma non ne abbiamo prova, non possiamo provare che dialettica sia struttura del reale \\
Non bisogna far diventare dialettica un dogma \t che pero lo è diventata \t non teme i fatti \\
Questo è il problema del marxismo \t davanti a qualsiasi esperienza, il marxista non cambia idea \t invece esperienza deve indurre in cambiamento se necessario \t non ci puo essere lettura aprioristica della realta \\
Distingue marxismo vivente e marxismo ufficiale (dogma, verita assoluta, etc.) \t mentre vivente legge la realta, è interpretativo, euristico, si confronta con la realta e cerca la verita \\
Il marxismo ufficiale è invece un marx. irrigidito dentro le sue certezze \t e questo perche ha generato con la dialettica un beato ottimismo, che afferma che si va sempre verso una sintesi migliore \t e questo serve per guidare la massa degli uomini \\
Questo marxismo hegeliano rende l uomo uno strumento passivo della dialettica \t e deresponsabilizza l uomo, in questa prospettiva uomo diventa pedina di realta che si muove cosi \\
Il marxista vero (vivente) ricerca nella realta, è euristico = è ricerca \\
Invece marx. uffi. è diventato un sapere eterno \t non sa + nulla, non vuole acquisire conoscenze, ma si vuole solo costituire a priori come un sapere assoluto \\
Questo marxismo ha dato vita all'esistenzialismo \t perche uomo nella dialettica si perde \t quindi esistienzialismo ha potuto nascere e affermare la realta indivudale dell uomo \\
Di fatto critica alla fine Hegel, perche critica dialettica
\end{document}
